% Diskussion — claude-tunnel bachelor thesis.
\chapter{Diskussion}
\label{ch:diskussion}

Dieses Kapitel bewertet das Ergebnis gegen die Forschungsfragen und das
Bedrohungsmodell (Kapitel~\ref{ch:anforderungen}) und benennt die offenen
Grenzen.

\section{Beantwortung der Forschungsfragen}
\label{sec:disk-ff}

\paragraph{FF1 (Realisierbarkeit).} Ein providerblinder Tunnel liess sich
konstruieren. Der Edge relayt nachweislich nur Chiffretext: der Ende-zu-Ende mit
Noise verschluesselte Verkehr wird byteweise durchgereicht, ohne dass der
Vermittler Ziel oder Nutzlast erhaelt. Rendezvous mit Proof-of-Work-Gating,
direkter Peer-to-Peer-Pfad mit Relay-Rueckfall und der
\gls{ac:tls}-ueber-\gls{ac:tcp}-Fallback funktionieren im Zusammenspiel; die
Kette Konto\,$\rightarrow$\,Aufladung\,$\rightarrow$\,guthabengedeckte
Token-Ausgabe\,$\rightarrow$\,Tunnel ist ebenfalls durchgaengig belegt. FF1 ist
damit fuer Architektur und Kernmechanismen positiv beantwortet.

\paragraph{FF2 (Latenz-Overhead).} Der Eigenanteil des Tunnels ist gering (rund
8\,ms Baseline) und der Overhead skaliert vorhersagbar linear mit der
Netzverzoegerung ($\text{RT} \approx 8{,}8 + 6{,}1 \cdot d$, siehe
Kapitel~\ref{ch:evaluation}). Der Faktor $\approx 6$ ist die wesentliche Kosten:
er ruehrt vom mehrfachen Queren der verzoegerten Strecke durch Handshake,
Rendezvous und Relay-Hin-/Rueckweg und liesse sich durch Verbindungs-Wiederverwendung
(0-RTT, Keep-Alive) deutlich senken.

\paragraph{FF3 (Netz- und Betriebsartverhalten).} Paketverlust trifft den Tail,
nicht den Median; die Betriebsart ist bei verlustfreiem Netz nicht
latenzbestimmend. Ein statistisch belastbarer Modus-Effekt unter Verlust liess
sich bei $n=20$ nicht nachweisen -- ein Hinweis, dass die Transportwahl fuer die
reine Latenz zweitrangig ist und primaer nach Semantik (Stream vs.\ Datagramm)
und Zensurresistenz getroffen werden sollte.

\section{Schutzziele gegen das Bedrohungsmodell}
\label{sec:disk-schutz}

Gegenueber dem neugierigen Betreiber (A1) werden Vertraulichkeit und
Ziel-Blindheit (S1, S2) strukturell erreicht: er sieht nur Chiffretext und opake
Tokens. Gegenueber dem aktiven Angreifer (A2) verhindert das Origin-Pinning die
Man-in-the-Middle-Uebernahme (S3); ein grober \gls{ac:udp}-Block wird durch den
\gls{ac:tcp}-Fallback ueberbrueckt. Gegen den Fluter (A3) daempft der
asymmetrische Proof-of-Work-Aufwand (S4). Nicht adressiert bleiben -- wie im
Bedrohungsmodell festgelegt -- Traffic-Analyse gegen A2 sowie die Oekonomie des
finanzierten Sybil-Angreifers (A4).

Ueber das Testbett hinaus erweitert die Produktivierung
(Kapitel~\ref{ch:produktivierung}) die Schutzziele um Identitaet und
Abrechnung: eine konventionelle OIDC-Authentifizierung bindet jede Aktion an
ein verifiziertes Subjekt, die Abrechnungsintegritaet haengt an einer
signierten Provider-Bestaetigung (der Kontrollknoten kann nicht selbst
gutschreiben), und ein Ratenlimit je Konto begrenzt die Token-Ausgabe. Die
strukturelle Blindheit des Betreibers bleibt davon unberuehrt -- die
zusaetzlichen Kontrollen betreffen Identitaet und Abrechnung, nicht den Zugriff
auf die Nutzlast.

\section{Offene Risiken und Grenzen}
\label{sec:disk-risiken}

\begin{description}
  \item[Finanzierter Sybil-Angriff (A4).] Die in der Produktivierung
    eingefuehrten konventionellen Konten verteuern Missbrauch, loesen die
    Sybil-Oekonomie aber nicht: wer zahlt, kann beliebig viele Tokens erwerben.
    Eine Abwehr muesste die Kosten pro schaedlicher Aktion ueber den Token-Preis
    hinaus anheben oder eine kontogebundene Reputation einfuehren.
  \item[Traffic-Analyse.] Groessen- und Timing-Seitenkanaele bleiben fuer einen
    beobachtenden Vermittler oder Netz-Angreifer sichtbar. Verschleierung (etwa
    Padding, konstante Raten) waere eine orthogonale Erweiterung mit eigenem
    Overhead.
  \item[NAT-Hole-Punching.] Der direkte Pfad gelingt im flachen Container-Testbett
    trivial; echtes gleichzeitiges Oeffnen ueber symmetrisches \gls{ac:nat} wurde
    nicht unter realen \gls{ac:nat}-Bedingungen erprobt und bleibt eine
    Testbett-Grenze.
  \item[PoW-Parametrisierung.] Die Schwierigkeit ist ein Kompromiss: zu niedrig
    schwaecht die Schranke, zu hoch benachteiligt schwache Clients. Eine
    adaptive, lastabhaengige Wahl wurde nicht untersucht.
\end{description}

\section{Methodische Einordnung}
\label{sec:disk-methodik}

Die Messung diente der Charakterisierung des Overheads, nicht einem Vergleich
mit Fremdsystemen. Das einseitige Impairment und der Frisch-Dial je Iteration
zeichnen ein konservatives Bild; ein groesseres $n$ und eine breitere Matrix
(inkl.\ Bandbreitenbegrenzung und PoW-Schwierigkeit) wuerden insbesondere die
Tail-Aussagen und einen etwaigen Modus-Effekt unter Verlust schaerfen.
